\begin{table}[H]
\centering
\caption{The representative point, retrained rather than re-measured. Every row re-runs the same learners from scratch on a state tensor built around a different anchor: identical seeds, identical record-level split, identical hyper-parameters, identical decoding grid, test scored once. Part I is the segmenter (3 seeds, 40 epochs) and is reported as mean $\pm$ sd over repeated runs, because the same input does not reproduce the same run; Part II is the frozen production detector (V8, 3 seeds, 18 epochs).}
\label{tab:p1-anchor}
\scriptsize
\begin{tabular}{lrrrrrr}
\toprule
Anchor & \multicolumn{3}{c}{Part I segmenter} & \multicolumn{3}{c}{Part II detector} \\
\cmidrule(lr){2-4}\cmidrule(lr){5-7}
 & ARI (mean $\pm$ sd) & runs & $\Delta$ARI (pp) & AP & AUC & Worst env. \\
\midrule
wrist & 0.5278 $\pm$ 0.0313 & 5 & +0.00 & 0.1493 & 0.8271 & 0.7777 \\
palm & 0.5440 $\pm$ 0.0355 & 5 & +1.62 & 0.1236 & 0.8075 & 0.7393 \\
tips & 0.5613 $\pm$ 0.0292 & 5 & +3.35 & 0.0819 & 0.8099 & 0.7519 \\
\bottomrule
\end{tabular}
\par\smallskip\footnotesize Wrist is the anchor the main pipeline uses, so its row is the reference. The two learners disagree: tips is the best anchor for the segmenter and the worst for the detector, so the choice of representative point is load-bearing and must be declared per part rather than fixed once for the paper. The pooled within-anchor sd is 0.0320 ARI; tips minus wrist is +3.35 points with a standard error of 1.92 points, which falls inside the noise.
\end{table}
